\documentclass[11pt]{article}
\usepackage[a4paper,margin=25mm]{geometry}
\usepackage{amsmath,amssymb,amsthm,booktabs,array,longtable}
\usepackage[T1]{fontenc}
\usepackage{lmodern}
\usepackage[hidelinks]{hyperref}
\newtheorem{theorem}{Theorem}
\newtheorem{proposition}[theorem]{Proposition}
\newtheorem{corollary}[theorem]{Corollary}
\hypersetup{pdftitle={Exact contents in Copeland's polynomial-content conjecture via Carlitz's congruence},pdfauthor={Haobo Ma and Wenlin Zhang}}
\title{Exact contents in Copeland's polynomial-content conjecture\\via Carlitz's congruence}
\author{Haobo Ma\\\small ChronoAI Pte Ltd, Singapore
\and Wenlin Zhang\\\small National University of Singapore\\\small ChronoAI Pte Ltd, Singapore}
\date{November 19, 2026}
\begin{document}
\maketitle
\begin{abstract}
For Copeland's polynomial-content conjecture (OEIS A131758, revision 48,
COMMENTS block signed September 9, 2008), we determine every row content
using Carlitz's polynomial congruence (1953/54, Theorem 3, equation (4.8)).
The content is $1$ at indices $0,1$ and all even indices, and $2$ at every
odd index at least $3$. This replaces the proposed divisor-count formula;
every index $6s+8$, $s\geq0$, also excludes its one-place shift.
The proof identifies content with a fixed divisor of an integer Laguerre
polynomial. We give direct proofs of associated transform identities and
correct normalization formulas in the linked OEIS entries.
\end{abstract}

\section{The assertion and the array}
In the block identified above, Copeland labels the following a conjecture:
``The greatest common divisor of the coefficients of each polynomial is
given by A060872(n)/n or, equivalently, by A038548(n)'' \cite{copeland}.
We replace this formula by an exact determination of the content.
The polynomial congruence used below is due to Carlitz \cite{carlitz};
Krylov and Li give a later direct proof \cite{krylov}. We include a
self-contained proof of its prime-modulus case and then determine the
fixed divisor exactly.

For nonnegative integers $n$ define the ordinary Laguerre polynomial and
its monic integer multiple by
\begin{equation}\label{laguerre}
 L_n(x)=\sum_{j=0}^n(-1)^j\binom nj\frac{x^j}{j!},\qquad
 B_n(x)=(-1)^n n!L_n(x)
 =\sum_{j=0}^n(-1)^{n+j}\binom nj\frac{n!}{j!}x^j.
\end{equation}
This is the parameter-zero normalization of DLMF 18.5.12 \cite{dlmf}.
Copeland's coefficient formula is
\begin{equation}\label{coefficient}
 a(n,m)=\sum_{k=0}^m(-1)^k\binom{n+1}{k}B_n(m-k+1),\qquad
 P_n(t)=\sum_{m=0}^n a(n,m)t^m \quad(0\leq m\leq n).
\end{equation}
Let $c_n$ denote the nonnegative gcd of the $n+1$ integer coefficients
of $P_n$. For a positive integer $u$ write
\[
 D(u)=\#\{d\in\mathbb Z_{>0}:d\mid u,\ d^2\leq u\}.
\]
Thus $D(u)=\mathrm{A038548}(u)=\mathrm{A060872}(u)/u$
\cite{divisors,products}. The conjecture is $c_n=D(n)$ for $n\geq1$.
The alternative $c_n=D(n+1)$ is not asserted in the source; it also fails.

\begin{proposition}\label{properties}
For every $n\geq0$, the polynomial $P_n$ has integer coefficients, degree
$n$, leading coefficient $n!$, and $P_n(1)=n!$. Moreover
\begin{equation}\label{series}
 P_n(t)=(1-t)^{n+1}\sum_{r\geq0}B_n(r+1)t^r,
 \qquad
 c_n=\gcd\{B_n(v):v\in\mathbb Z_{>0}\}.
\end{equation}
\end{proposition}
\begin{proof}
Every coefficient in \eqref{laguerre} is integral since $j!\mid n!$.
Let $\Delta f(x)=f(x+1)-f(x)$. Since $B_n$ is monic of degree $n$,
$\Delta^{n+1}B_n=0$ and $\Delta^n B_n=n!$.
For $m\leq n$ the coefficient of $t^m$ in the series on the right of
\eqref{series} is \eqref{coefficient}; for $m>n$ it is
$\Delta^{n+1}B_n(m-n)=0$. This proves the series identity.
The relation $\Delta^{n+1}B_n(0)=0$ gives
$a(n,n)=(-1)^n B_n(0)=n!$.
Also the coefficient of $t^n$ in $P_n(t)/(1-t)$ is both the row sum and
\[
 \sum_{k=0}^n(-1)^k\binom nk B_n(n+1-k)=\Delta^n B_n(1)=n!.
\]
Finally, multiplication by $(1-t)^{n+1}$ and its integer-coefficient
inverse preserves the ideal generated by all coefficients of a power
series. Applied to \eqref{series}, this identifies the two gcds.
All statements include $P_0=1$.
\end{proof}

\section{The exact replacement formula}
\begin{theorem}\label{content}
The contents of Copeland's array are
\begin{equation}\label{answer}
 c_n=\begin{cases}
 1,& n=0,\ n=1,\ \text{or }n\text{ is even},\\
 2,& n\geq3\text{ is odd}.
 \end{cases}
\end{equation}
\end{theorem}
\begin{proof}
We first exclude every odd prime from the fixed divisor in \eqref{series}.
For a prime $p$, write $n=qp+r$, $0\leq r<p$. Terms of
\eqref{laguerre} with $j<qp$ contain a factor $p$ in $n!/j!$.
For $j=qp+s$, $0\leq s\leq r$, coefficient comparison in
$(1+z)^n\equiv(1+z^p)^q(1+z)^r\pmod p$ gives
$\binom n{qp+s}\equiv\binom rs$, and $n!/(qp+s)!\equiv r!/s!$.
Consequently
\begin{equation}\label{primecongruence}
 B_n(x)\equiv x^{qp}B_r(x)\pmod p.
\end{equation}
This is the ordinary-Laguerre, prime-modulus specialization of Carlitz's
congruence \cite[Theorem 3, (4.8)]{carlitz}.
Suppose $p$ is odd. If $r<p-1$, the monic polynomial $B_r$ has fewer
than $p-1$ roots, so it is nonzero at some nonzero residue.
If $r=p-1$ and it vanished at every nonzero residue, it would equal
$x^{p-1}-1$ over $\mathbb F_p$. This is impossible: the coefficient of
$x^{p-2}$ in $B_{p-1}$ is $-(p-1)^2\equiv-1$, whereas that in
$x^{p-1}-1$ is zero. In either case \eqref{primecongruence} supplies
a positive integer $v$ with $p\nmid B_n(v)$, excluding $p\mid c_n$.

For $n\geq2$, reduction modulo $2$ gives $B_n(x)\equiv x^n$ if $n$
is even, and $B_n(x)\equiv x^{n-1}(x+1)$ if $n$ is odd.
Indeed, terms with $j\leq n-2$ have even factorial quotient, and the
next coefficient is $-n^2$. Thus even $n$ gives odd $B_n(1)$ and
$c_n=1$, while odd $n\geq3$ gives even values at every integer and
$2\mid c_n$.

To exclude a factor $4$ when $n\geq3$ is odd, terms with $j\leq n-4$
vanish modulo $4$. The top four coefficients of $B_n$ are
\[
 1,\quad -n^2,\quad \frac{n^2(n-1)^2}{2},\quad
 -\frac{n^2(n-1)^2(n-2)^2}{6}.
\]
The last two are both $0$ modulo $4$ if $n\equiv1\pmod4$, and both
$2$ if $n\equiv3\pmod4$; division by $3$ is legitimate modulo $4$.
At $x=3$ their combined contribution is zero modulo $4$, since the
corresponding powers are $3$ and $1$ modulo $4$. The first two
contribute $3-1=2$. Therefore $B_n(3)\equiv2\pmod4$ and $c_n=2$.
The remaining cases are $P_0=1$ and $P_1=t$.
\end{proof}

\begin{corollary}\label{equalities}
For positive $n$, the source equation $c_n=D(n)$ holds exactly at
$n=1,2$ and at odd integers $n=p^2,p^3,pq$, where $p,q$ are distinct
odd primes in the last expression.
The unposed shifted equation $c_n=D(n+1)$ holds exactly at $n=0,1$,
at even $n\geq2$ with $n+1$ prime, and at odd indices
$n=3,7,2q-1$ with $q$ an odd prime.
In particular, both equations fail at every $n=6s+8$, $s\geq0$.
\end{corollary}
\begin{proof}
Divisor pairing gives $D(u)=\lceil\tau(u)/2\rceil$, where $\tau(u)$ is
the number of positive divisors. Thus $D(u)=1$ exactly for $u=1$ or
prime $u$, and $D(u)=2$ exactly for $u=p^2,p^3,pq$ with distinct primes
in the last case. Substitute \eqref{answer}; an even integer in the
second list is $4$, $8$, or $2q$ with $q$ odd prime.
Finally $6s+8$ is even and composite, and $6s+9$ is composite.
Both divisor counts exceed its content $1$.
\end{proof}

For comparison, Greubel's existing table \cite{greubel}, flattened entries
$36$--$44$, gives
\begin{equation}\label{row8}
 (a(8,m))_{m=0}^8=(6209,-41417,48645,376259,-1724861,
 3233925,-2977609,1078849,40320).
\end{equation}
The two-coefficient identity
$-6417\cdot6209-962\cdot(-41417)=1$ already proves $c_8=1$,
whereas $D(8)=D(9)=2$. The larger family in Corollary~\ref{equalities}
contains $n=2^{12s+7}$ for every $s\geq0$.

\begin{corollary}\label{parity}
For even $n\geq2$, $P_n(t)\equiv(1+t)^{n-1}\pmod2$.
For odd $n\geq3$, $P_n(t)\equiv0\pmod2$.
\end{corollary}
\begin{proof}
For even $n$, $B_n(r+1)\equiv r+1\pmod2$, whose generating series is
$1/(1-t)^2$. Use \eqref{series}. For odd $n\geq3$, all integer
values of $B_n$ are even, as in the theorem.
\end{proof}
In A009940, $Q_n=n!L_n(1)$, so the correct first-column identification
is $a(n,0)=(-1)^n Q_n$, not $Q_n$ without the sign: for example the
two values at $n=3$ are $4$ and $-4$.
Bala's congruence $Q_{n+k}\equiv(-1)^k Q_n\pmod k$ \cite{constantcolumn}
already implies its constant-column parity ($k=2$).
It follows from the corollary on page 5 of Bala's 2017 note \cite{bala},
taking $F(x)=G(x)=1/(1-x)$ there. Carlitz's earlier polynomial
congruence supplies the whole-row parity in Corollary~\ref{parity}.

\section{Sibling identities and normalization corrections}
All generating-series identities below are formal in $x$ with rational
functions of $t$ as coefficients; polynomial identities extend across
removable specializations. Set
$R_n(t)=tP_n(t)/(n!(1-t)^{n+1})$ and $G(x,t)=\sum_{n\geq0}P_n(t)x^n/n!$.
\begin{corollary}\label{transforms}
Writing $y=x(t-1)$, one has
\begin{align}
 \sum_{n\geq0}R_n(t)z^n
 &=\frac{t}{(1+z)(\exp(-z/(1+z))-t)},\label{rgf}\\
 G(x,t)&=\frac{t-1}{(1-y)(t-\exp(y/(1-y)))}.\label{egf}
\end{align}
Let $s_j=\mathrm{A000262}(j)$, characterized by
$\sum s_j u^j/j!=\exp(u/(1-u))$. Then the modified Riordan recurrence
stated in A131758 is
\begin{equation}\label{mra}
 P_n(t)=n!(t-1)^n+\sum_{j=0}^{n-1}\binom nj s_{n-j}
 (t-1)^{n-j-1}P_j(t)\quad(n\geq1).
\end{equation}
If $E_n(t)$ denotes the Eulerian polynomial normalized by
$\sum_{v\geq1}v^n t^v=tE_n(t)/(1-t)^{n+1}$, including $E_0=1$, then
\begin{equation}\label{eulerian}
 \frac{E_n(t)}{n!}=\sum_{j=0}^n\binom nj(1-t)^{n-j}\frac{P_j(t)}{j!}.
\end{equation}
The inverse replaces $E_n$ on the left by $P_n$, and $P_j$ on the right
by $E_j$ and $(1-t)^{n-j}$ by $(t-1)^{n-j}$. In particular the
alternating row sum is
\[
 P_n(-1)=n!\sum_{j=0}^n\binom nj(-2)^{n-j}\frac{E_j(-1)}{j!}.
\]
\end{corollary}
\begin{proof}
Summing the finite expansion of $L_n$ gives
$\sum L_n(v)(-z)^n=(1+z)^{-1}\exp(vz/(1+z))$.
Sum this against $t^v$, $v\geq1$, to obtain \eqref{rgf}, and substitute
$z=x(1-t)$ to obtain \eqref{egf}.
Multiplying \eqref{egf} by $t-\exp(y/(1-y))$ and extracting coefficients
gives \eqref{mra}. Finally \eqref{laguerre} yields
\[
 R_n(t)=\sum_{j=0}^n(-1)^{n-j}\binom nj
 \frac{\operatorname{Li}_{-j}(t)}{j!},\qquad
 \operatorname{Li}_{-j}(t)=\sum_{v\geq1}v^j t^v.
\]
Binomial inversion and the definition of $E_n$ give \eqref{eulerian}
and its inverse. These are the corresponding A131758 transform formulas.
\end{proof}

\begin{corollary}\label{related}
Let $A_n(t)$ be the row polynomials of A131202, including $A_0=1$,
with the generating function specified in that entry \cite{related}.
Then
\begin{equation}\label{Aidentity}
 A_n(t)=P_n(t)-n(t-1)P_{n-1}(t)\quad(n\geq1),\qquad
 P_n(t)=\sum_{k=0}^n\binom nk A_{n-k}(t)k!(t-1)^k.
\end{equation}
Consequently the multiplier sequence in its reconstruction of $P_n$
is $k!(t-1)^k$, rather than the printed
``c(n,t) = n!*(t-1)\^{}(n-1) for n > 0''.
At $n=1$ the printed rule gives $2$, whereas $P_1(t)=t$.
At $t=2$ the identity is $A_n(2)=P_n(2)-nP_{n-1}(2)$, as stated in
A084358. The binomial matrix with entries
$\binom nk A_{n-k}(2)$ is A133289 and its inverse is $2I-$A129652.
Also $A_n(1)=n!$, and $A_n(t)$ has degree $n-1$ and leading
coefficient $s_n$ for $n\geq1$, as stated in A131202.
\end{corollary}
\begin{proof}
The A131202 generating function is $(t-1)/(t-\exp(y/(1-y)))$.
By \eqref{egf} it equals $(1-y)G(x,t)$, proving both formulas in
\eqref{Aidentity}. At $t=2$ it becomes $1/(2-\exp(x/(1-x)))$, the
A084358 generating function \cite{lists}. Its reciprocal is
$2-\exp(x/(1-x))$. Multiplication of exponential generating functions
is binomial convolution, so their binomial matrices are inverses.
The latter reciprocal has constant coefficient $1$ and coefficients
$-s_n$ for $n>0$, exactly $2I-$A129652 \cite{riordan,inverse}.
At $t=1$, \eqref{Aidentity} and Proposition~\ref{properties} give
$A_n(1)=n!$. To identify the leading coefficient, replace $x$ by $x/t$
in the generating function of $A_n$. Coefficientwise as $t\to\infty$,
\[
 \sum_{n\geq0}\frac{A_n(t)}{t^n}\frac{x^n}{n!}
 =1+\frac{\exp(x/(1-x))-1}{t}+O(t^{-2}).
\]
The two terms of degree $n$ in \eqref{Aidentity} cancel, so $A_n$ has
degree at most $n-1$; the displayed limit identifies its leading
coefficient as the positive integer $s_n$.
\end{proof}

\begin{corollary}\label{fubini}
Let $F_n=\mathrm{A000670}(n)$ be the ordered Bell numbers, defined by
$\sum F_n z^n/n!=1/(2-e^z)$, and put $b_n=(-1)^nF_n$.
The reciprocal Laguerre-transform pair involving $P_n(2)$ in A131758
holds with $b_n$ as the signed ordered Bell sequence.
The printed identity ``OB(n) = -Li(-n,2)'' instead requires
\begin{equation}\label{half}
 b_n=-\tfrac12\operatorname{Li}_{-n}(2)\quad(n\geq0).
\end{equation}
In the same specialization, $P_n(2)/n!=(-1)^{n+1}R_n(2)/2$.
\end{corollary}
\begin{proof}
For fixed $\alpha$, define the generalized Laguerre transform on a
sequence $u_j$ by
$ (\mathcal T_\alpha u)_n=\sum_{j=0}^n(-1)^j(n!/j!)
\binom{n+\alpha}{n-j}u_j$.
If $f(z)=\sum u_jz^j/j!$, its transformed generating function is
$(1-z)^{-\alpha-1}f(-z/(1-z))$. The substitution
$z\mapsto-z/(1-z)$ is an involution and its two prefactors cancel,
so $\mathcal T_\alpha^2$ is the identity. This also proves the
generalized reciprocal identity stated in A131758.
For $\alpha=0$, apply this formula to $f(z)=1/(2-e^{-z})$.
The result is \eqref{egf} at $t=2$, proving the reciprocal pair.
Moreover
\[
 \sum_{n\geq0}-\operatorname{Li}_{-n}(2)\frac{z^n}{n!}
 =\frac{2e^z}{2e^z-1}=\frac{2}{2-e^{-z}},
\]
which proves \eqref{half}; these are rational continuations of the
polylogarithms from $|t|<1$. The final formula follows directly from
the definition of $R_n$.
\end{proof}

\section{Coverage and verification}
\small
\begin{longtable}{@{}>{\raggedright\arraybackslash}p{0.52\linewidth}>{\raggedright\arraybackslash}p{0.43\linewidth}@{}}
\toprule Source item and status & Coverage\\\midrule\endhead
A131758 content conjecture $c_n=D(n)$ & Refuted-and-corrected by Theorem~\ref{content}.\\
Shift $c_n=D(n+1)$ & Not asserted in the source; fails at $n=8$ and all $6s+8$. Exact equality cases: Corollary~\ref{equalities}.\\
Integer rows, leading coefficient and row sum & Proved in Proposition~\ref{properties}.\\
First column identified with A009940 & Refuted-and-corrected sign identification; explicit values $4,-4$ at $n=3$.\\
A009940 constant-column congruence & Prior proof: Carlitz (1953/54), then Krylov--Li (2014/2016); Bala's general proof (2017) specializes to his OEIS formula (2023). Whole-row consequence: Corollary~\ref{parity}.\\
A131758 rational generating function, MRA recurrence, binomial/Eulerian transforms & Proved in Corollary~\ref{transforms}.\\
A131202 reconstruction, leading coefficient and row sum & Refuted-and-corrected reconstruction exponent; leading coefficient and row sum proved in Corollary~\ref{related}. Other listed identifications are not covered.\\
A084358 finite-difference identity; A133289 inverse matrix & Proved in Corollary~\ref{related}, with the earlier source formulas cited.\\
A131758 reciprocal Laguerre transforms and signed ordered Bell normalization & Proved, and factor-of-two formulas refuted-and-corrected, in Corollary~\ref{fubini}.\\
Other operational/Bessel/zeta formulas, complex interpolation and the 2017 Todd-operator comparison & Not covered.\\
\bottomrule
\end{longtable}
\normalsize
The \href{https://github.com/the-omega-institute/trureturing/blob/742ce086eb9c8e6e36cde0fcdf712f63260c045c/D5/S0/Certificates/CopelandA131758PolynomialContentRefutation.lean}{pinned Lean file}
defines \texttt{a}, \texttt{P}, \texttt{IsSourceRow}, \texttt{rowContent},
\texttt{D}, and \texttt{claim} in the namespace
\texttt{D5.S0.Certificates.CopelandA131758PolynomialContentRefutation}.
Its theorem \texttt{result : $\neg$ claim} negates the joint assertion
of integer-row existence for every $n\geq0$ and content $D(n)$ for every
positive-index source row. Its argument uses the two coefficients
$6209,-41417$ and $D(8)=2$. The exact-content formula, global integrality,
equality classifications, parity law and sibling identities have the
written proofs above. The accompanying \texttt{verify.py} supplies finite
exact-arithmetic checks of those statements; its scope and results are in
\texttt{README.md} and \texttt{verification.json}.

\begin{thebibliography}{14}
\bibitem{copeland} T.~Copeland, \emph{Coefficients of numerators of rational
functions whose binomial transforms give the normalized polylogarithms
Li(-n,t)/n!}, OEIS A131758, revision 48, May 23, 2024; COMMENTS block
signed September 9, 2008. \url{https://oeis.org/A131758}.
\bibitem{dlmf} NIST, \emph{Digital Library of Mathematical Functions},
18.5.12, specialized to $\alpha=0$. \url{https://dlmf.nist.gov/18.5.E12}.
\bibitem{divisors} OEIS Foundation Inc., \emph{Number of divisors of n that
are at most sqrt(n)}, A038548, revision 195, September 24, 2026.
\url{https://oeis.org/A038548}.
\bibitem{products} OEIS Foundation Inc., \emph{Sum of d*d' over all unordered
pairs (d,d') with d*d' = n}, A060872, revision 41, October 21, 2024;
formula credited to O.~E.~Pol, May 5, 2020. \url{https://oeis.org/A060872}.
\bibitem{greubel} G.~C.~Greubel, \emph{Table of n, a(n) for the first 50 rows,
flattened}, A131758, entries 36--44.
\url{https://oeis.org/A131758/b131758.txt}.
\bibitem{constantcolumn} OEIS Foundation Inc., \emph{a(n) = n!*L\_n(1)},
A009940, revision 63, July 18, 2025; congruence credited to P.~Bala,
March 12, 2023. \url{https://oeis.org/A009940}.
\bibitem{related} T.~Copeland, \emph{A coefficient tree from the list
partition transform}, A131202, revision 28, November 1, 2026.
\url{https://oeis.org/A131202}.
\bibitem{lists} OEIS Foundation Inc., \emph{Lists of sets of lists}, A084358,
revision 40, September 8, 2022; generating function credited to V.~Jovovic,
September 28, 2003; finite-difference identity to T.~Copeland,
September 30, 2007. \url{https://oeis.org/A084358}.
\bibitem{riordan} P.~Barry, \emph{Exponential Riordan array
[e\^{}(x/(1-x)),x]}, A129652, revision 42, March 13, 2026; binomial
formula credited to P.~Bala, May 14, 2012. \url{https://oeis.org/A129652}.
\bibitem{inverse} T.~Copeland, \emph{Riordan matrix T from A084358},
A133289, revision 19, September 10, 2015. \url{https://oeis.org/A133289}.
\bibitem{fubini} OEIS Foundation Inc., \emph{Fubini numbers: number of
ordered partitions}, A000670. \url{https://oeis.org/A000670}.
\bibitem{carlitz} L.~Carlitz, \emph{Congruence properties of the polynomials
of Hermite, Laguerre and Legendre}, Mathematische Zeitschrift 59
(1953/54), 474--483, especially p.~479, Theorem~3, (4.8).
\url{https://doi.org/10.1007/BF01180274}.
\bibitem{krylov} N.~A.~Krylov and Z.~Li, \emph{A congruence property of
irreducible Laguerre polynomials in two variables}, Integers 16 (2016),
A57, Section~3, Corollary~2, (10), pp.~7--8; earlier preprint
arXiv:1408.1904v1 (August 8, 2014), Corollary~3, (16).
\url{https://math.colgate.edu/~integers/q57/q57.pdf}.
\bibitem{bala} P.~Bala, \emph{Integer sequences that become periodic on
reduction modulo k for all k}, December 3, 2017, p.~5, Corollary.
\url{https://oeis.org/A047974/a047974_1.pdf}.
\end{thebibliography}
\noindent\small\textcopyright\ 2026 Haobo Ma and Wenlin Zhang.
Licensed under Creative Commons Attribution 4.0 International (CC BY 4.0),
\url{https://creativecommons.org/licenses/by/4.0/}.
\end{document}
